\section{Introduction}

Multilingual legal retrieval requires more than ordinary cross-lingual text matching. Legal users search for obligations, exceptions, applicability limits, evidentiary requirements, and consequences across language-specific realizations of the same act, while benchmark supervision must remain auditable and stable enough to support retrieval evaluation. The JRC-Acquis corpus is a natural starting point for this setting because it provides broad multilingual coverage of European Union legal and regulatory texts together with a stable document identifier, \texttt{CELEX}, that links language versions of the same act.

Our goal is to turn that raw multilingual legal collection into a reproducible retrieval benchmark rather than only a cleaned corpus dump. This requires an end-to-end pipeline that starts with XML parsing and text normalization, separates retrieval-oriented document views from generation-oriented prompt views, selects source documents that are substantive enough to support good legal questions, and then constructs multilingual relevance labels in a way that remains transparent about why documents are linked. We therefore build supervision from same-language legal question--answer generation on selected target-language document realizations and define positive relevance at the retained sampled \texttt{CELEX}-group level.

The resulting benchmark is designed for practical retrieval evaluation. It exports standard retrieval artifacts for Hugging Face and MTEB-style use, preserves a richer construction trace for auditing, and explicitly addresses the question-shape failures that are common in legal QA generation, such as provision-led wording, inventory-style prompts, and coordinated multi-part questions. The remainder of this document summarizes the benchmark contributions, the full construction method, the evaluation protocol, and the current limitations of the approach.
